\chapter{Микроскопический гамильтониан повторных взаимодействий}
\label{ch:v8-microscopic-repeated-interaction-hamiltonian-gate}

\section{Звёздное взаимодействие на минимальной среде}

Пусть $D_a=D_a^*$, $a=1,\ldots,12$, --- полный вещественный
межсекторных стрелок базис на системном носителе $\mathcal H_S=\mathbb C^{21}$, а
\begin{equation}
 G=\sum_{a=1}^{12}D_a^2.
\end{equation}
Минимальная Стайнспринг-среда предыдущих глав имеет вид
\begin{equation}
 \mathcal K_{\rm env}
 =\mathbb C|0\rangle\oplus
 \operatorname{span}_{\mathbb C}\{|a\rangle:a=1,\ldots,12\}.
\end{equation}
На $\mathcal H_S\otimes\mathcal K_{\rm env}\cong\mathbb C^{273}$ определим
\begin{equation}
 H_{\rm int}
 =\sum_{a=1}^{12}D_a\otimes
 \bigl(|a\rangle\langle0|+|0\rangle\langle a|\bigr).
 \label{eq:v8-star-interaction-hamiltonian}
\end{equation}
Это не логарифм произвольного Стайнспринг-унитарный: оператор построен прямо
из уже существующих стрелок и вакуумной линии среды. Точная LCF-проверка даёт
\begin{equation}
 H_{\rm int}=H_{\rm int}^*,
 \qquad
 \langle0|H_{\rm int}^2|0\rangle=G.
 \label{eq:v8-star-interaction-second-moment}
\end{equation}

\section{Предел слабых столкновений}

Для длительности микрошагa $h$ положим
\begin{equation}
 U_h=\exp\bigl(-i\sqrt h\,H_{\rm int}\bigr)
\end{equation}
и подготовим вспомогательную систему в $|0\rangle$. Соответствующие блоки Крауса имеют
разложения
\begin{align}
 K_0(h)&=I-\frac h2G+O(h^2),\\
 K_a(h)&=-i\sqrt h\,D_a+O(h^{3/2}).
\end{align}
Поэтому редуцированная карта удовлетворяет
\begin{equation}
 \Psi_h(X)=X+h\mathcal L_{q\ell}(X)+O(h^2),
 \qquad
 \mathcal L_{q\ell}(X)
 =\sum_aD_aXD_a-\frac12\{G,X\}.
 \label{eq:v8-star-interaction-gksl-tangent}
\end{equation}
Повторение со свежими вспомогательными системами и масштабом $h=u/n$ возвращает уже
сертифицированный предел
\begin{equation}
 \lim_{n\to\infty}\Psi_{u/n}^{\,n}=e^{u\mathcal L_{q\ell}}.
\end{equation}
Тем самым предъявлен явный унитарный микромеханизм безразмерной
марковской динамики.

\section{Калибровочная ковариантность и пространство связей}

При калибровочном преобразовании репер $(D_a)$ смешивается вещественной
ортогональной матрицей. Если скачковая часть среды несёт двойственное
представление, свёртка \eqref{eq:v8-star-interaction-hamiltonian}
инвариантна и не зависит от выбора ортонормированного базиса.

Однако эта симметрия не выбирает единичную свёртку. Точная система
коммутантных уравнений для двенадцатимерного репер даёт
\begin{equation}
 \dim_{\mathbb R}\operatorname{Comm}_G(E_{q\ell}^{\rm cross})=8,
 \qquad
 \dim_{\mathbb R}\operatorname{Comm}^{\rm sym}_G
 (E_{q\ell}^{\rm cross})=4.
 \label{eq:v8-star-coupling-commutant-dimensions}
\end{equation}
Причина состоит в том, что $QLYR$ и $XLdR$ являются двумя эквивалентными
вещественными калибровочными копиями. Их полный коммутант имеет тип $M_2(\mathbb C)$,
и каждый вещественный элемент этого восьмимерного коммутанта задаёт
самосопряжённый $H_{\rm int}$. Редуцированный генератор зависит от
положительной симметричной метрики $C^TC$; пространство симметричных
коммутантных метрик имеет размерность $4$. Диагональные скорости двух
семейств являются лишь двумерным подсемейством; калибровочная симметрия допускает
также межсемейное смешивание и микроскопические фазы, невидимые одной
метрике скоростей.

\section{Почему точный конечный канал не воспроизведён}

Звёздная экспонента и ранее выбранный точный канал
\begin{equation}
 \Phi_h(X)=\sqrt{I-hG}\,X\sqrt{I-hG}
 +h\sum_aD_aXD_a
\end{equation}
имеют одинаковую производную в нуле, но разные коэффициенты порядка $h^2$.
Если $J(X)=\sum_aD_aXD_a$, то
\begin{align}
 [h^2]\Psi_h(X)
 &=\frac14GXG+\frac1{24}\{G^2,X\}
   -\frac16\{G,J(X)\},\\
 [h^2]\Phi_h(X)
 &=\frac14GXG-\frac18\{G^2,X\}.
\end{align}
Точный матричный свидетель различия уже возникает на первом диагональном
матричных единиц. Следовательно, $H_{\rm int}$ выводит требуемый генератор и
непрерывный предел, но не задним числом объясняет конкретную нелинейную
форму конечного Краус-шагa.

Кроме того, замена $H_{\rm int}\mapsto gH_{\rm int}$ умножает скорость
генератора на $g^2$. Поэтому абсолютный перевод $u$ в физическое время и
подготовка последовательности свежих вспомогательных систем остаются внешними данными.

\begin{theorem}[Микроскопическая реализация межсекторных стрелок полугруппы]
Генератор межсекторных стрелок Тома VIII допускает явную самосопряжённую,
калибровочно-ковариантную реализацию повторных взаимодействий на минимальной среде
размерности $13$. Её предел слабых столкновений равен
$e^{u\mathcal L_{q\ell}}$. Но калибровочно допустимое пространство вещественных
самосопряжённых связей взаимодействия имеет размерность $8$, а пространство
симметричных скоростных метрик --- размерность $4$;
конечный унитарный шаг отличается от выбранной точной Краус-карты начиная с
порядка $h^2$.
\end{theorem}

\begin{nogocriterion}[Запрет объявления единственного гамильтониан часов]
Нельзя считать свёртку с единичной матрицей связей выведенной одной
калибровочной, вещественно-структурной и эрмитовой типизацией. Нельзя отождествлять существование
микроскопический гамильтониан столкновения с выводом абсолютной длительности такта,
автономной поставки свежей среды или точного конечного канала $\Phi_h$.
Для дальнейшей селекции требуется дополнительная геометрическая структура,
различающая или канонически смешивающая две эквивалентные межсекторные семьи.
\end{nogocriterion}

\begin{protocol}
\begin{enumerate}
 \item системная размерность: \textbf{$21$};
 \item размерность минимальной среды: \textbf{$13$};
 \item полный унитарный носитель: \textbf{$273$};
 \item самосопряжённость $H_{\rm int}$: \textbf{точно};
 \item второй момент вакуума: \textbf{$G$};
 \item GKSL-касательная: \textbf{точно};
 \item предел слабых столкновений: \textbf{получен};
 \item калибровочная ковариантность: \textbf{точно};
 \item самосопряжённые калибровочные связи: \textbf{$8$ вещественных измерений};
 \item симметричные скоростные метрики $C^TC$: \textbf{$4$ вещественных измерения};
 \item единственный гамильтониан взаимодействия: \textbf{нет};
 \item совпадение с точным конечным $\Phi_h$: \textbf{нет, различие $O(h^2)$};
 \item абсолютная скорость: \textbf{не выведена};
 \item автономная свежая среда: \textbf{не выведена};
 \item следующий гейт: \textbf{геометрический селектор матрицы связей}.
\end{enumerate}
\end{protocol}

Машинный сертификат сохранён в
\path{s2t_v8_microscopic_repeated_interaction_hamiltonian_gate_results.json}.

\begin{thebibliography}{9}
\bibitem{v8-micro-attal-pautrat}
S.~Attal, Y.~Pautrat,
\emph{From repeated to continuous quantum interactions},
Annales Henri Poincar\'e \textbf{7} (2006), 59--104;
arXiv:math-ph/0311002.
\bibitem{v8-micro-attal-joye}
S.~Attal, A.~Joye,
\emph{Weak coupling and continuous limits for repeated quantum interactions},
Journal of Statistical Physics \textbf{126} (2007), 1241--1283;
arXiv:math-ph/0501012.
\bibitem{v8-micro-erker}
P.~Erker et al.,
\emph{Autonomous quantum clocks: does thermodynamics limit our ability to
measure time?}, Physical Review X \textbf{7} (2017), 031022;
arXiv:1609.06704.
\end{thebibliography}